\section{Discussion and Conclusion}\label{sec:discussion}
Sparse visual codes and order-theoretic descriptions play different
roles. The surrogate supplies candidate coordinates; the context states
exactly how those coordinates are shared and retrieved. The construction
retains the original relation-induced proof pattern and is fully
compatible with established pattern structures and ordinally scaled
FCA. Its value here is an explicit, checkable specification of visual
feature queries rather than a new general representation theorem.

The main vision-specific lesson is that max pooling changes the spatial
question. A coordinatewise image query allows separate witnesses for its
requirements. A same-site query requires a common witness and generally
cannot be represented by a single summary vector. The single-vector
criterion and the downset context explain this distinction exactly,
while the executed pilot shows that it materially affects the retrieved
sets even in a small convolutional network.

The empirical findings also constrain the utility claim. Graded queries
outperform the chosen presence baseline, but cosine retrieval is stronger
for digit labels. Such labels measure class association, not the
correctness of a semantic name attached to a latent coordinate. Replacing
all activations by SAE reconstructions checks one aspect of fidelity;
it does not establish the causal importance of an individual coordinate
or conjunction. Controlled feature interventions, as explored in related
vision SAE work~\cite{stevens2025vision}, address a different question.

Several limitations remain. The digit benchmark uses one small dataset and one
fixed split.
The natural-image extension uses a deterministic CIFAR-10 subset,
one frozen ResNet34, and a surrogate with only moderate reconstruction
fidelity. It supplies inspectable examples rather than a second
comparative retrieval benchmark.
Training seeds and overlapping source draws do not provide independent
dataset replications. Training and SAE seeds are varied together, so
these runs do not isolate sources of variation. Dead coordinates are
common, there is no external spatial annotation, and no vision
transformer or transcoder is evaluated. The spatial downset construction is
analyzed
mathematically but its full lattice is not mined experimentally. These
limitations preclude broad claims about vision-model interpretability.

A larger study should fix pretrained backbone and surrogate revisions,
verify their activation conventions, and compare at least two vision
families on independent natural-image collections. It should include
spatially annotated tasks, matched sparsity and reconstruction quality,
strong dense and sparse retrieval baselines, and uncertainty estimates
whose resampling units match the intended population. Coordinate
alignment across independently learned dictionaries requires its own
method; coincident index numbers do not supply such alignment. A useful
intervention study would compare selected feature edits with random and
norm-matched controls while measuring reconstruction error and the
backbone's output change.

The present paper provides a self-contained mathematical specification,
an executable digit pilot, measured natural-image feature galleries, and an
explicit route to this broader study.
It shows how to ask precise questions about shared sparse evidence, and
why retaining the spatial meaning of those questions is essential.
